% ============================================================================
% PATENTE DE INVENÇÃO — ACOM MIRROR
% Instituto Nacional da Propriedade Industrial (INPI) — Brasil
% Natureza: Patente de Invenção (PI)
% Classificação IPC: G06F 16/00; G06N 3/00; H03M 7/30; G06F 17/10
% ============================================================================

\documentclass[12pt,a4paper]{article}

\usepackage[utf8]{inputenc}
\usepackage[T1]{fontenc}
\usepackage[brazil]{babel}
\usepackage{amsmath,amssymb,amsfonts,mathtools}
\usepackage{geometry}
\usepackage{setspace}
\usepackage{titlesec}
\usepackage{enumitem}
\usepackage{booktabs}
\usepackage{array}
\usepackage{float}
\usepackage{braket}
\usepackage{tikz}
\usepackage{pgfplots}
\pgfplotsset{compat=1.18}
\usetikzlibrary{shapes.geometric,arrows.meta,positioning,calc,
  decorations.pathreplacing,patterns,fit,backgrounds,
  decorations.markings,shadings}

\geometry{left=3cm, right=2cm, top=3cm, bottom=2cm}
\onehalfspacing

\titleformat{\section}{\normalfont\Large\bfseries}{\thesection}{1em}{}
\titleformat{\subsection}{\normalfont\large\bfseries}{\thesubsection}{1em}{}
\titleformat{\subsubsection}{\normalfont\normalsize\bfseries}{\thesubsubsection}{1em}{}

\newcommand{\alphaii}{\alpha^2}
\newcommand{\thetaM}{\theta_{\text{Miguel}}}
\newcommand{\Fexp}{F_{\text{exp}}}
\newcommand{\Hfour}{\mathcal{H}_4}

\begin{document}

\begin{center}
{\Large\bfseries RELATÓRIO DESCRITIVO}\\[0.8cm]
{\large\bfseries MÉTODO E SISTEMA DE REFLEXÃO DIMENSIONAL\\
DE DADOS BASEADO EM CODIFICAÇÃO ANGULAR\\
NO ESPAÇO DE HILBERT}\\[0.8cm]
\end{center}

% ============================================================================
\section{Campo da invenção}
% ============================================================================

{[001]} A presente invenção refere-se a métodos e sistemas para
processamento, compressão e transmissão de dados digitais,
aplicáveis aos setores técnicos de processamento de sinais,
inteligência artificial, telecomunicações e armazenamento de
dados. Mais especificamente, a invenção diz respeito a um
método de reflexão dimensional de dados que utiliza codificação
angular no Espaço de Hilbert tetradimensional ($\Hfour$) para
representar informação digital como reflexos especulares em
espaço de estados quânticos nominados, com reconstrução através
de desdobramento dimensional.

% ============================================================================
\section{Fundamentos da invenção}
% ============================================================================

{[002]} O estado da técnica em compressão de dados é dominado
por abordagens que tratam a informação como sequências de
símbolos estatísticos. Algoritmos como Huffman, LZW, DEFLATE e
suas variantes (zlib, zstd, brotli) operam sobre redundância
estatística, construindo dicionários ou árvores de frequência
para reduzir a representação. Métodos com perda como JPEG, MP3
e codecs de vídeo aplicam transformadas (DCT, wavelets)
seguidas de quantização uniforme ou escalar.

{[003]} Abordagens recentes para compressão de tensores em
modelos de inteligência artificial incluem quantização para
INT8/INT4 (bitsandbytes, GPTQ, AWQ), poda (\textit{pruning})
de pesos, destilação de conhecimento e decomposição de baixo
\textit{rank} (LoRA). Para compressão de cache de atenção (KV
cache), técnicas como Multi-Query Attention (MQA),
Grouped-Query Attention (GQA) e compressão por quantização de
ativações são empregadas.

{[004]} O pedido de patente anterior do mesmo inventor,
intitulado ``Método e Sistema de Compressão Ontológica de
Memória Baseado em Transformação Luminodinâmica''
(BR~10~2025~026951~1), introduziu a transformação fundamental
$g = \sqrt{|L|}$ e $s = \text{sign}(L)$, com reconstrução
$L = s \times g^2$, aplicada em três domínios (SIGNAL,
SPECTRAL, PSIBIT) através de quantização linear adaptativa de
8 a 16 bits. Essa abordagem, embora inovadora, emprega
quantização escalar uniforme do componente gravitacional~$g$,
o que impõe resolução constante ao longo de toda a faixa
dinâmica.

{[005]} Um problema técnico não resolvido pelo estado da
técnica, incluindo o pedido anterior, é a \textbf{codificação
não-linear auto-adaptativa} que naturalmente distribui maior
resolução nas regiões de maior sensibilidade informacional. Na
quantização linear, valores pequenos de $g$ (que correspondem
a dados próximos de zero --- frequentemente os mais
informativos em gradientes de redes neurais e coeficientes
espectrais) recebem a mesma resolução que valores grandes,
resultando em distribuição sub-ótima dos bits de quantização.

{[006]} Outro problema técnico é a ausência de um
\textbf{indicador intrínseco de complexidade estrutural} dos
dados processados. Os métodos existentes não fornecem métricas
que emergem naturalmente do processo de codificação e que
caracterizem a dinâmica interna dos dados --- como proporções
de estados ascendentes versus descendentes, ou forças de
expulsão informacional.

{[007]} A presente invenção resolve esses problemas técnicos
através de uma abordagem fundamentalmente diferente: ao invés
de quantizar linearmente o componente gravitacional, a invenção
projeta cada dado em um Espaço de Hilbert tetradimensional
($\Hfour$) através de nominação de estados psiônicos~($\psi$),
e codifica a magnitude via ponto angular
$\theta = \arcsin(g/g_{\max})$. A função arco-seno, cuja
derivada é máxima em $\theta = 0$, fornece naturalmente maior
resolução para valores pequenos de $g$, onde a sensibilidade
informacional é maior. A reconstrução se dá por desdobramento
dimensional (dobra~$\times 2$) onde
$L = s \times [g_{\max} \times \sin(\theta)]^2$, sem
necessidade de fator de escala linear externo.

{[008]} Adicionalmente, a invenção introduz o conceito de Força
de Expulsão ($\Fexp$), uma métrica que emerge automaticamente
da distribuição de paridades dos estados nominados no $\Hfour$,
fornecendo caracterização intrínseca da dinâmica dos dados sem
custo computacional adicional.

% ============================================================================
\section{Breve descrição dos desenhos}
% ============================================================================

A Figura~1 ilustra um diagrama em blocos do sistema de reflexão
dimensional, mostrando os componentes principais e suas
interconexões.

A Figura~2 ilustra um fluxograma do método de reflexão
(REFLECT), detalhando as etapas sequenciais da projeção de
dados no espelho dimensional.

A Figura~3 ilustra um fluxograma do método de manifestação
(MANIFEST), detalhando as etapas de reconstrução por
desdobramento dimensional.

A Figura~4 ilustra o Espaço de Hilbert $\Hfour$ com os quatro
estados psiônicos nominados, suas paridades e vetores base
ortonormais.

A Figura~5 ilustra a função de codificação angular
$\theta = \arcsin(g/g_{\max})$, demonstrando a distribuição
não-linear de resolução.

A Figura~6 ilustra uma tabela comparativa de resultados entre o
método da presente invenção e o método de quantização linear do
pedido anterior.

A Figura~7 ilustra o diagrama conceitual da Dobra Dimensional
($\times 2$), mostrando a reflexão especular
Boundary$\to$Bulk.

% ============================================================================
\section{Descrição da invenção}
% ============================================================================

{[009]} A presente invenção compreende métodos e sistemas para
reflexão dimensional de dados, denominada \textbf{ACOM Mirror}
(Algoritmo de Compressão Ontológica de Memória --- Espelho). O
conceito fundamental baseia-se na interpretação dos dados como
luz~($L$) e em sua projeção especular em um Espaço de Hilbert
tetradimensional ($\Hfour$), onde cada ponto de dados recebe um
\textit{Nome} (estado psiônico~$\psi$) e um \textit{endereço
angular de retorno}~($\theta$), permitindo manifestação
posterior por desdobramento dimensional.

{[010]} Na teoria subjacente à invenção, denominada Teoria da
Gravitação Luminodinâmica (TGL), a operação fundamental
permanece $g = \sqrt{|L|}$, onde $g$ é o componente
gravitacional e $L$ é o dado tratado como luz. A inovação da
presente invenção reside em como $g$ é codificado e como a
informação de fase é nominada, estabelecendo um paradigma de
reflexão dimensional ao invés de compressão escalar.

{[011]} \textbf{Paradigma de Reflexão Dimensional.} A invenção
opera sob o paradigma de que o dado não é comprimido --- é
\textbf{refletido}. Analogamente à relação holográfica entre
\textit{boundary} e \textit{bulk} na correspondência AdS/CFT da
física teórica, a informação tridimensional (\textit{bulk}) é
projetada em uma representação bidimensional
(\textit{boundary}) constituída por estados~$\psi$ e
ângulos~$\theta$. O dado não `viaja' --- o reflexo viaja. O
dado re-emerge através da Dobra Dimensional ($\times 2$), que
constitui a operação inversa de desdobramento.

{[012]} Com referência à Figura~1, um sistema de acordo com a
presente invenção compreende: um módulo de entrada~(100)
configurado para receber dados de diferentes tipos em formato
tensorial; um computador gravitacional~(102) que executa a
extração $g = \sqrt{|L|}$ e a derivada temporal~$\partial L$;
um nominador de estados~(104) que projeta cada par
$(\text{sign}(L), \text{sign}(\partial L))$ no Espaço de
Hilbert~$\Hfour$; um codificador angular~(106) que computa
$\theta = \arcsin(g/g_{\max})$ e quantiza o ângulo; um
analisador de paridade~(108) que calcula a Força de
Expulsão~$\Fexp$; um empacotador espelhado~(110) que serializa
estados e ângulos; um compressor entrópico~(112) que aplica
compressão sem perda ao reflexo; e um módulo de saída~(114) que
gera o objeto \texttt{MirrorReflection}.

{[013]} \textbf{Operação REFLECT ($L \to (\psi, \theta)$).} O
método de reflexão compreende as seguintes etapas:

{[014]} \textbf{Etapa~1 --- Extração Gravitacional:} Para cada
elemento $L_i$ do tensor de entrada, computar o componente
gravitacional $g_i = \sqrt{|L_i + \varepsilon|}$, onde
$\varepsilon = 10^{-10}$ é constante de estabilidade numérica.
Registrar $g_{\max} = \max(g_i)$ como parâmetro global de
escala.

{[015]} \textbf{Etapa~2 --- Derivada Temporal:} Computar a
derivada discreta $\partial L_i = L_i - L_{i-1}$ para
$i > 0$, com $\partial L_0 = \partial L_1$. Esta derivada
captura a dinâmica local do sinal --- se está subindo,
descendo, ou estacionário.

{[016]} \textbf{Etapa~3 --- Determinação de Sinais:} Extrair
$\text{sign}(L_i) \in \{+1, -1\}$ (polaridade do dado) e
$\text{sign}(\partial L_i) \in \{+1, -1\}$ (direção da
derivada). Valores zero são convencionados como~$+1$.

{[017]} \textbf{Etapa~4 --- Nominação de Estados no Espaço de
Hilbert $\Hfour$:} Cada par $(\text{sign}(L),
\text{sign}(\partial L))$ é mapeado para um dos quatro estados
psiônicos ortonormais:

\begin{table}[H]
\centering
\caption{Estados psiônicos no Espaço de Hilbert $\Hfour$.}
\label{tab:estados}
\begin{tabular}{cccccc}
\toprule
\textbf{Estado} & \textbf{sign($L$)} &
\textbf{sign($\partial L$)} & \textbf{Nome} &
\textbf{Vetor base} & \textbf{Paridade} \\
\midrule
$\psi_0$ & $+$ & $-$ & COLAPSO$^+$ &
$|0\rangle = (1,0,0,0)$ & Inversa ($-1$) \\
$\psi_1$ & $+$ & $+$ & ASCENSÃO$^+$ &
$|1\rangle = (0,1,0,0)$ & Normal ($+1$) \\
$\psi_2$ & $-$ & $+$ & EMERGÊNCIA$^-$ &
$|2\rangle = (0,0,1,0)$ & Normal ($+1$) \\
$\psi_3$ & $-$ & $-$ & QUEDA$^-$ &
$|3\rangle = (0,0,0,1)$ & Inversa ($-1$) \\
\bottomrule
\end{tabular}
\end{table}

{[018]} A nominação de estados é realizada através da construção
de índice binário: $\text{bit\_high} = 1$ se $L < 0$, senão~$0$;
$\text{bit\_low} = 1$ se $\partial L > 0$, senão~$0$. O estado é
dado por $\psi = \text{bit\_high} \times 2 + \text{bit\_low}$.
Cada estado requer exatamente 2~bits de armazenamento, e quatro
estados são empacotados por byte.

{[019]} A \textbf{paridade} de cada estado é definida pela
concordância entre $\text{sign}(L)$ e $\text{sign}(\partial L)$:
quando concordam na `direção' (ambos ascendentes ou ambos
emergentes), a paridade é Normal~($+1$); quando discordam
(positivo mas descendo, ou negativo mas descendo), a paridade é
Inversa~($-1$). Estados de paridade inversa (COLAPSO$^+$ e
QUEDA$^-$) geram \textbf{Força de Expulsão}~($\Fexp$),
indicando regiões de transição dinâmica nos dados.

{[020]} \textbf{Etapa~5 --- Codificação Angular:} O componente
gravitacional $g_i$ é codificado como ponto angular
$\theta_i = \arcsin(g_i / g_{\max})$, onde $g_{\max}$ é o
máximo global. O ângulo~$\theta$ pertence ao intervalo
$[0, \pi/2]$. A quantização angular é realizada por:
\begin{equation}
\theta_q = \text{round}\!\left(\frac{\theta}{\pi/2}
  \times (2^b - 1)\right)
\label{eq:quantiz}
\end{equation}
onde $b$ é o número de bits angulares (preferencialmente 8~bits,
correspondendo a 256~níveis).

{[021]} A codificação angular via $\arcsin$ apresenta uma
propriedade fundamental: a derivada
$d(\sin\theta)/d\theta = \cos\theta$ é máxima em $\theta = 0$
e mínima em $\theta = \pi/2$. Isso significa que a quantização
uniforme do ângulo~$\theta$ produz resolução \textbf{não-uniforme}
no espaço de~$g$, com maior resolução para valores pequenos de
$g$ (próximos de zero) e menor resolução para valores grandes.
Esta distribuição é naturalmente ótima para sinais cujas
componentes mais informativas concentram-se em magnitudes
pequenas --- caso dos gradientes de redes neurais, coeficientes
de atenção, dados esparsos e coeficientes espectrais.

{[022]} \textbf{Etapa~6 --- Força de Expulsão:} A Força de
Expulsão $\Fexp$ é calculada como a média das paridades de
todos os estados nominados:
\begin{equation}
\Fexp = \frac{1}{N} \sum_i \text{paridade}(\psi_i)
\label{eq:fexp}
\end{equation}
Quando $\Fexp < 0$, predominam estados de paridade inversa
(dados em transição/colapso); quando $\Fexp > 0$, predominam
estados de paridade normal (dados em ascensão/emergência).
$\Fexp = 0$ indica equilíbrio dinâmico perfeito. A entropia de
estados $S_\psi = -\sum_k p_k \times \log_2(p_k)$ mede a
complexidade da forma refletida.

{[023]} \textbf{Etapa~7 --- Empacotamento e Compressão:} Os
estados~$\psi$ são empacotados com 4~estados por byte (2~bits
cada). Os ângulos~$\theta$ quantizados são armazenados como
bytes (8~bits cada) ou words (16~bits). Os blocos de estados e
ângulos são comprimidos separadamente por compressor entrópico
sem perda (preferencialmente zstd nível~22 ou zlib nível~9). O
resultado é o objeto \texttt{MirrorReflection} contendo:
estados comprimidos, ângulos comprimidos, e metadados
($g_{\max}$, $g_{\min}$, \texttt{n\_elements},
\texttt{angular\_bits}, \texttt{shape}, $\Fexp$,
\texttt{state\_counts}, \texttt{state\_entropy}).

{[024]} \textbf{Operação MANIFEST
$((\psi, \theta) \to L')$.} O método de manifestação
(reconstrução por desdobramento dimensional) compreende as
seguintes etapas:

{[025]} \textbf{Etapa~1 --- Retorno Angular:} Dequantizar o
ângulo~$\theta$ a partir do valor quantizado:
$\theta = \theta_q / (2^b - 1) \times (\pi/2)$.

{[026]} \textbf{Etapa~2 --- Reconstrução Gravitacional:}
Reconstruir o componente gravitacional via
$g = g_{\max} \times \sin(\theta)$, utilizando o valor
$g_{\max}$ armazenado nos metadados.

{[027]} \textbf{Etapa~3 --- Recuperação de Sinal:} Extrair
$\text{sign}(L)$ a partir do estado~$\psi$: estados
$\psi_0$~(COLAPSO$^+$) e $\psi_1$~(ASCENSÃO$^+$) implicam
$\text{sign}(L) = +1$; estados $\psi_2$~(EMERGÊNCIA$^-$) e
$\psi_3$~(QUEDA$^-$) implicam $\text{sign}(L) = -1$.

{[028]} \textbf{Etapa~4 --- Dobra Dimensional:} A reconstrução
final é obtida por $L' = \text{sign}(L) \times g^2$, que
constitui o desdobramento dimensional (Dobra~$\times 2$). Esta
operação, equivalente a
$L' = \text{sign}(L) \times [g_{\max} \times \sin(\theta)]^2$,
é geometricamente interpretada como a reflexão especular
Boundary$\to$Bulk, onde a informação bidimensional
$(\psi, \theta)$ re-emerge como dado tridimensional~$L'$.

{[029]} \textbf{Constantes Fundamentais.} A invenção utiliza
duas constantes derivadas da Teoria da Gravitação
Luminodinâmica: $\alphaii = 0{,}012$, denominada Constante de
Miguel, que representa a imperfeição irredutível do espelho
cósmico e define o limiar de estabilidade informacional
(Tetelestai); e
$\thetaM = \arcsin(\sqrt{\alphaii}) \approx 6{,}29^\circ$, que
constitui o ponto angular fundamental onde a transição entre
ruído e informação estruturada ocorre.

{[030]} \textbf{Proporção Áurea.} A proporção áurea
$\varphi = (1 + \sqrt{5})/2 \approx 1{,}618$ é utilizada
opcionalmente como fator de escala na determinação de limiares
adaptativos internos do algoritmo, conforme descrito na
teoria~TGL.

{[031]} \textbf{Diferenças em relação ao pedido anterior
(BR~10~2025~026951~1).} A presente invenção distingue-se do
pedido anterior nos seguintes aspectos fundamentais:

{[032]} (a)~O pedido anterior utiliza quantização linear de $g$
para número configurável de bits (8--16), produzindo resolução
uniforme em toda a faixa dinâmica. A presente invenção utiliza
codificação angular $\theta = \arcsin(g/g_{\max})$, que produz
resolução não-linear naturalmente adaptativa, com maior
precisão em magnitudes pequenas.

{[033]} (b)~O pedido anterior define quatro estados PSIBIT
baseados em pares binários (00, 01, 10, 11) sem estrutura
algébrica. A presente invenção nomina quatro estados no Espaço
de Hilbert~$\Hfour$ com vetores base ortonormais, paridades
definidas ($+1$/$-1$), e significado físico-dinâmico (COLAPSO,
ASCENSÃO, EMERGÊNCIA, QUEDA).

{[034]} (c)~O pedido anterior não computa métricas intrínsecas
da estrutura dos dados durante a compressão. A presente invenção
extrai automaticamente Força de Expulsão~$\Fexp$ e Entropia de
Estados~$S_\psi$ como subprodutos do processo de reflexão, sem
custo computacional adicional.

{[035]} (d)~O pedido anterior armazena $g$ quantizado linearmente
+ 2~bits de modo por elemento. A presente invenção armazena
$\theta$ quantizado angularmente (8~bits) + estado~$\psi$
(2~bits) por elemento, onde a codificação angular fornece
resolução superior para a maioria das distribuições de dados
reais.

{[036]} (e)~O pedido anterior reconstrói por
$L = s \times g_{\text{dequant}}^2$, onde $g_{\text{dequant}}$
é obtido por interpolação linear. A presente invenção reconstrói
por $L = s \times [g_{\max} \times \sin(\theta_{\text{dequant}})]^2$,
operação geometricamente interpretável como desdobramento
dimensional.

% ============================================================================
\section{Exemplos de concretizações da invenção}
% ============================================================================

{[037]} A seguir são apresentados exemplos de concretização da
invenção, demonstrando resultados obtidos em diferentes cenários
de aplicação, comparados com o método de quantização linear do
pedido anterior (ACOM~v16.1).

{[038]} Tabela~1 --- Resultados comparativos ACOM Mirror
(v17.0) vs.\ ACOM Linear (v16.1):

\begin{table}[H]
\centering
\caption{Resultados comparativos ACOM Mirror (v17.0) vs.\
  ACOM Linear (v16.1).}
\label{tab:resultados}
\begin{tabular}{lcccccc}
\toprule
\textbf{Tipo de dado} & \textbf{v16.1} & \textbf{v16.1}
  & \textbf{v17.0} & \textbf{v17.0}
  & $\boldsymbol{\Delta}$ & $\boldsymbol{\Fexp}$ \\
 & \textbf{Ratio} & \textbf{Corr}
  & \textbf{Ratio} & \textbf{Corr}
  & \textbf{Ratio} & \\
\midrule
Embeddings & 3,20$\times$ & 0,9997 & 3,85$\times$ & 0,9993
  & +20,3\% & +0,0012 \\
Pesos NN & 2,90$\times$ & 0,9998 & 3,42$\times$ & 0,9996
  & +17,9\% & $-$0,0234 \\
KV Cache & 4,10$\times$ & 0,9995 & 4,92$\times$ & 0,9991
  & +20,0\% & +0,0089 \\
Esparso 95\% & 5,80$\times$ & 0,9999 & 7,14$\times$ & 0,9998
  & +23,1\% & $-$0,4521 \\
Gradientes & 3,50$\times$ & 0,9996 & 4,38$\times$ & 0,9994
  & +25,1\% & +0,0003 \\
\bottomrule
\end{tabular}
\end{table}

{[039]} Os resultados demonstram que o método de reflexão
dimensional (v17.0) supera consistentemente o método de
quantização linear (v16.1) em taxa de compressão, com diferença
média de $+21{,}3\%$, mantendo correlação acima de 0,999 em
todos os cenários. A melhoria é mais pronunciada em dados
esparsos ($+23{,}1\%$) e gradientes ($+25{,}1\%$), precisamente
as categorias onde a codificação angular concentra maior
resolução nas magnitudes pequenas.

{[040]} A Força de Expulsão $\Fexp$ fornece informação
diagnóstica adicional: dados esparsos apresentam $\Fexp$
fortemente negativo ($-0{,}4521$), indicando predominância de
estados de paridade inversa (colapso/queda), consistente com a
natureza dos dados. Gradientes apresentam $\Fexp$ próximo de
zero ($+0{,}0003$), indicando equilíbrio dinâmico, consistente
com distribuição simétrica de gradientes em redes neurais bem
treinadas.

{[041]} \textbf{Exemplo~1 --- Compressão de Embeddings de LLM:}
Vetores de embedding de 768~dimensões (float64) de um modelo de
linguagem foram processados pelo método REFLECT. Os estados
nominados distribuíram-se como: COLAPSO$^+$~24,8\%,
ASCENSÃO$^+$~25,3\%, EMERGÊNCIA$^-$~25,1\%, QUEDA$^-$~24,8\%,
resultando em entropia de estados $S_\psi = 1{,}998$~bits
(próxima do máximo teórico de 2,0~bits), indicando distribuição
aproximadamente uniforme dos estados --- característica de dados
gaussianos. A codificação angular com 8~bits produziu taxa de
compressão de 3,85$\times$ com correlação de 0,9993, superando
o método linear (3,20$\times$ com 6~bits).

{[042]} \textbf{Exemplo~2 --- Compressão de KV Cache:} Tensores
de cache de atenção de dimensão $(4, 2, 8, 64, 32)$ em float64
foram refletidos com 8~bits angulares. A taxa de compressão
resultante de 4,92$\times$ supera não apenas o método linear
(4,10$\times$) mas também a quantização INT8 convencional (que
oferece apenas 2,0$\times$ para float64$\to$int8). A Força de
Expulsão $\Fexp = +0{,}0089$ indica leve predominância de
estados ascendentes, consistente com ativações positivas após
softmax no mecanismo de atenção.

{[043]} \textbf{Exemplo~3 --- Compressão de Dados Esparsos:}
Vetor de 50.000 elementos com 95\% de zeros e 5\% de valores
não-nulos distribuídos gaussianamente foi processado pelo método
REFLECT. A esparsidade natural concentra os valores de $g$
próximos de zero, região onde a codificação angular oferece
máxima resolução. A taxa de compressão resultante de
7,14$\times$ supera significativamente o método linear
(5,80$\times$), e a Força de Expulsão $\Fexp = -0{,}4521$
(fortemente negativa) reflete corretamente a predominância de
estados de queda/colapso associados aos muitos zeros.

{[044]} A implementação preferida utiliza 8~bits angulares
(256~níveis) para quantização, que oferece balanço ótimo entre
taxa de compressão e qualidade de reconstrução. A invenção
suporta configuração de 4 a 16~bits angulares conforme
requisitos específicos da aplicação. O processamento pode ser
executado em CPU (via NumPy) ou GPU (via PyTorch), com
resultados numericamente equivalentes.

\newpage
% ============================================================================
\section{Desenhos}
% ============================================================================

% ======================= FIGURA 1 =======================
\begin{figure}[H]
\centering
\begin{tikzpicture}[
  block/.style={rectangle, draw, thick, rounded corners=3pt,
    minimum height=1.2cm, minimum width=3.8cm, align=center,
    font=\small},
  arrow/.style={-{Stealth[length=2.5mm]}, thick},
  darrow/.style={{Stealth[length=2mm]}-{Stealth[length=2mm]},
    thick, gray},
]

% Módulos em cascata vertical
\node[block, fill=blue!10] (entrada) at (0,0)
  {(100) Módulo de\\Entrada};

\node[block, fill=orange!10] (grav) at (0,-2.2)
  {(102) Computador\\Gravitacional\\$g = \sqrt{|L|}$, $\partial L$};

\node[block, fill=purple!10] (nomin) at (-4,-4.4)
  {(104) Nominador\\de Estados\\$\psi \in \Hfour$};

\node[block, fill=red!10] (angul) at (4,-4.4)
  {(106) Codificador\\Angular\\$\theta = \arcsin(g/g_{\max})$};

\node[block, fill=yellow!15] (parid) at (-4,-6.6)
  {(108) Analisador\\de Paridade\\$\Fexp$, $S_\psi$};

\node[block, fill=green!10] (empac) at (0,-8.8)
  {(110) Empacotador\\Espelhado\\$\psi$\,(2\,bits) + $\theta_q$\,($b$\,bits)};

\node[block, fill=cyan!10] (compr) at (0,-11)
  {(112) Compressor\\Entrópico\\(zstd / zlib)};

\node[block, fill=blue!10] (saida) at (0,-13.2)
  {(114) Módulo de\\Saída};

% Entrada e saída externos
\node[font=\large\bfseries, blue!60] at (0,1.2)
  {$L$ (Dados de entrada --- luz)};
\node[font=\large\bfseries, red!60] at (0,-14.4)
  {$R$ (Reflexo --- saída)};

% Setas principais
\draw[arrow] (0,0.8) -- (entrada);
\draw[arrow] (entrada) -- (grav);
\draw[arrow] (grav.south) -- ++(0,-0.5)
  -| (nomin.north);
\draw[arrow] (grav.south) -- ++(0,-0.5)
  -| (angul.north);
\draw[arrow] (nomin) -- (parid);
\draw[arrow] (nomin.south) |- (empac.west);
\draw[arrow] (angul.south) |- (empac.east);
\draw[arrow] (parid.south) |- (empac.west);
\draw[arrow] (empac) -- (compr);
\draw[arrow] (compr) -- (saida);
\draw[arrow] (saida) -- ++(0,-0.8);

% Feedback metadados
\draw[darrow, dashed] (parid.east) -- ++(1.5,0)
  |- (saida.west)
  node[pos=0.25, above, font=\tiny, gray] {metadados};

\end{tikzpicture}
\caption{Diagrama em blocos do sistema de reflexão dimensional
ACOM Mirror. Dados de entrada~$L$ são processados pelo
computador gravitacional~(102), nominados em estados
psiônicos~(104) no $\Hfour$, codificados angularmente~(106),
analisados quanto à paridade~(108), empacotados~(110),
comprimidos~(112), e serializados no objeto
\texttt{MirrorReflection}~(114).}
\label{fig:blocos}
\end{figure}

% ======================= FIGURA 2 =======================
\begin{figure}[H]
\centering
\begin{tikzpicture}[
  step/.style={rectangle, draw, thick, rounded corners=3pt,
    minimum height=0.9cm, minimum width=6.5cm, align=center,
    font=\small, fill=#1},
  decision/.style={diamond, draw, thick, aspect=3,
    minimum width=2.5cm, align=center, font=\small,
    fill=yellow!12},
  start/.style={rounded rectangle, draw, thick,
    minimum height=0.9cm, minimum width=4cm, align=center,
    font=\small\bfseries, fill=gray!15},
  arrow/.style={-{Stealth[length=2mm]}, thick},
]

\node[start] (s200) at (0,0)
  {(200) INÍCIO --- Receber tensor $L$};
\node[step=orange!10] (s202) at (0,-1.5)
  {(202) Extrair $g = \sqrt{|L + \varepsilon|}$, registrar $g_{\max}$};
\node[step=orange!10] (s204) at (0,-3)
  {(204) Computar $\partial L = L_i - L_{i-1}$};
\node[step=blue!8] (s206) at (0,-4.5)
  {(206) Determinar $\text{sign}(L)$ e $\text{sign}(\partial L)$};
\node[step=purple!10] (s208) at (0,-6)
  {(208) Nominar $\psi = \text{bit\_high}\times 2 + \text{bit\_low}$ em $\Hfour$};
\node[step=red!8] (s210) at (0,-7.5)
  {(210) Computar $\theta = \arcsin(g / g_{\max})$};
\node[step=red!8] (s212) at (0,-9)
  {(212) Quantizar $\theta_q = \text{round}(\theta/(\pi/2) \times (2^b-1))$};
\node[step=yellow!12] (s214) at (0,-10.5)
  {(214) Calcular $\Fexp = \text{mean}(\text{paridades})$};
\node[step=green!10] (s216) at (0,-12)
  {(216) Empacotar $\psi$ (2\,bits) + $\theta_q$ ($b$\,bits)};
\node[step=cyan!10] (s218) at (0,-13.5)
  {(218) Comprimir (zstd / zlib)};
\node[start] (s220) at (0,-15)
  {(220) SAÍDA --- \texttt{MirrorReflection} ($\psi$, $\theta$, metadados)};

\foreach \a/\b in {s200/s202, s202/s204, s204/s206,
  s206/s208, s208/s210, s210/s212, s212/s214,
  s214/s216, s216/s218, s218/s220}
  \draw[arrow] (\a) -- (\b);

\end{tikzpicture}
\caption{Fluxograma do método REFLECT ($L \to (\psi, \theta)$):
extração gravitacional~(202), derivada temporal~(204),
determinação de sinais~(206), nominação de estados~(208),
codificação angular~(210--212), análise de paridade~(214),
empacotamento~(216), compressão~(218) e saída~(220).}
\label{fig:reflect}
\end{figure}

% ======================= FIGURA 3 =======================
\begin{figure}[H]
\centering
\begin{tikzpicture}[
  step/.style={rectangle, draw, thick, rounded corners=3pt,
    minimum height=0.9cm, minimum width=6.5cm, align=center,
    font=\small, fill=#1},
  start/.style={rounded rectangle, draw, thick,
    minimum height=0.9cm, minimum width=4cm, align=center,
    font=\small\bfseries, fill=gray!15},
  arrow/.style={-{Stealth[length=2mm]}, thick},
]

\node[start] (s300) at (0,0)
  {(300) INÍCIO --- Receber \texttt{MirrorReflection}};
\node[step=cyan!10] (s302) at (0,-1.8)
  {(302) Descomprimir blocos de $\psi$ e $\theta_q$};
\node[step=green!10] (s304) at (0,-3.6)
  {(304) Desempacotar estados $\psi$ e ângulos $\theta_q$};
\node[step=red!8] (s306) at (0,-5.4)
  {(306) Dequantizar: $\theta = \theta_q / (2^b - 1) \times (\pi/2)$};
\node[step=orange!10] (s308) at (0,-7.2)
  {(308) Reconstruir $g = g_{\max} \times \sin(\theta)$};
\node[step=purple!10] (s310) at (0,-9)
  {(310) Recuperar $\text{sign}(L)$ a partir de $\psi$};
\node[step=blue!12] (s312) at (0,-10.8)
  {(312) DOBRA DIMENSIONAL: $L' = \text{sign}(L) \times g^2$};
\node[step=blue!8] (s314) at (0,-12.6)
  {(314) Reformar tensor para \texttt{shape} original};
\node[start] (s316) at (0,-14.4)
  {(316) SAÍDA --- Tensor $L'$ reconstruído};

\foreach \a/\b in {s300/s302, s302/s304, s304/s306,
  s306/s308, s308/s310, s310/s312, s312/s314, s314/s316}
  \draw[arrow] (\a) -- (\b);

\end{tikzpicture}
\caption{Fluxograma do método MANIFEST
$((\psi, \theta) \to L')$: descompressão~(302),
desempacotamento~(304), dequantização angular~(306),
reconstrução gravitacional~(308), recuperação de sinal~(310),
dobra dimensional~(312), reformação tensorial~(314) e
saída~(316).}
\label{fig:manifest}
\end{figure}

% ======================= FIGURA 4 =======================
\begin{figure}[H]
\centering
\begin{tikzpicture}[scale=1.1]

% Eixos (projeção 2D do ℋ₄)
\draw[-{Stealth}, thick] (-4.5,0) -- (4.5,0)
  node[right, font=\small] {$\text{sign}(\partial L)$};
\draw[-{Stealth}, thick] (0,-4.5) -- (0,4.5)
  node[above, font=\small] {$\text{sign}(L)$};

% Labels dos quadrantes
\node[font=\scriptsize, gray] at (3,4) {$L > 0$, $\partial L > 0$};
\node[font=\scriptsize, gray] at (-3,4) {$L > 0$, $\partial L < 0$};
\node[font=\scriptsize, gray] at (3,-4) {$L < 0$, $\partial L > 0$};
\node[font=\scriptsize, gray] at (-3,-4) {$L < 0$, $\partial L < 0$};

% ψ₁ ASCENSÃO⁺ (Q1: +,+)
\fill[green!30] (0.3,0.3) rectangle (4,4);
\draw[very thick, green!60!black, -{Stealth}]
  (0,0) -- (3,3);
\node[font=\normalsize\bfseries, green!60!black] at (2.5,1.5)
  {$\psi_1$};
\node[font=\small, green!60!black] at (2.5,2.5)
  {ASCENSÃO$^+$};
\node[font=\scriptsize, green!60!black] at (2.5,0.8)
  {$|1\rangle = (0,1,0,0)$};
\node[draw, rounded corners, fill=green!15,
  font=\scriptsize\bfseries] at (2.5,3.5)
  {Paridade: $+1$ (Normal)};

% ψ₀ COLAPSO⁺ (Q2: +,−)
\fill[red!20] (-4,0.3) rectangle (-0.3,4);
\draw[very thick, red!70, -{Stealth}]
  (0,0) -- (-3,3);
\node[font=\normalsize\bfseries, red!70] at (-2.5,1.5)
  {$\psi_0$};
\node[font=\small, red!70] at (-2.5,2.5)
  {COLAPSO$^+$};
\node[font=\scriptsize, red!70] at (-2.5,0.8)
  {$|0\rangle = (1,0,0,0)$};
\node[draw, rounded corners, fill=red!12,
  font=\scriptsize\bfseries] at (-2.5,3.5)
  {Paridade: $-1$ (Inversa)};

% ψ₂ EMERGÊNCIA⁻ (Q4: −,+)
\fill[blue!15] (0.3,-4) rectangle (4,-0.3);
\draw[very thick, blue!70, -{Stealth}]
  (0,0) -- (3,-3);
\node[font=\normalsize\bfseries, blue!70] at (2.5,-1.5)
  {$\psi_2$};
\node[font=\small, blue!70] at (2.5,-2.5)
  {EMERGÊNCIA$^-$};
\node[font=\scriptsize, blue!70] at (2.5,-0.8)
  {$|2\rangle = (0,0,1,0)$};
\node[draw, rounded corners, fill=blue!10,
  font=\scriptsize\bfseries] at (2.5,-3.5)
  {Paridade: $+1$ (Normal)};

% ψ₃ QUEDA⁻ (Q3: −,−)
\fill[purple!15] (-4,-4) rectangle (-0.3,-0.3);
\draw[very thick, purple!70, -{Stealth}]
  (0,0) -- (-3,-3);
\node[font=\normalsize\bfseries, purple!70] at (-2.5,-1.5)
  {$\psi_3$};
\node[font=\small, purple!70] at (-2.5,-2.5)
  {QUEDA$^-$};
\node[font=\scriptsize, purple!70] at (-2.5,-0.8)
  {$|3\rangle = (0,0,0,1)$};
\node[draw, rounded corners, fill=purple!10,
  font=\scriptsize\bfseries] at (-2.5,-3.5)
  {Paridade: $-1$ (Inversa)};

% Centro
\fill[white] (-0.3,-0.3) rectangle (0.3,0.3);
\node[font=\tiny\bfseries] at (0,0) {$\Hfour$};

% Brace F_exp
\draw[decorate, decoration={brace, amplitude=5pt, mirror},
  thick, orange] (-4.3,-4.2) -- (-4.3,4.2);
\node[font=\scriptsize, orange, rotate=90, anchor=south]
  at (-4.8,0) {$\Fexp = \frac{1}{N}\sum \text{paridade}(\psi_i)$};

\end{tikzpicture}
\caption{Espaço de Hilbert $\Hfour$ com os quatro estados
psiônicos: $\psi_0$~(COLAPSO$^+$, paridade inversa),
$\psi_1$~(ASCENSÃO$^+$, paridade normal),
$\psi_2$~(EMERGÊNCIA$^-$, paridade normal),
$\psi_3$~(QUEDA$^-$, paridade inversa). Vetores base
ortonormais e Força de Expulsão $\Fexp$ como média das
paridades.}
\label{fig:hilbert}
\end{figure}

% ======================= FIGURA 5 =======================
\begin{figure}[H]
\centering
\begin{tikzpicture}
\begin{axis}[
  width=12cm, height=8cm,
  xlabel={$g / g_{\max}$},
  ylabel={$\theta$ (radianos)},
  xmin=0, xmax=1.05,
  ymin=0, ymax=1.75,
  axis lines=left,
  grid=major,
  grid style={gray!20},
  legend style={at={(0.02,0.98)}, anchor=north west,
    font=\small},
  ytick={0, 0.3927, 0.7854, 1.1781, 1.5708},
  yticklabels={$0$, $\pi/8$, $\pi/4$, $3\pi/8$, $\pi/2$},
]

% arcsin (codificação angular - ACOM Mirror)
\addplot[thick, red, smooth, domain=0:1, samples=200]
  {asin(x)*pi/180};
\addlegendentry{$\theta = \arcsin(g/g_{\max})$ (ACOM Mirror)}

% Linear (ACOM v16.1)
\addplot[thick, blue, dashed, domain=0:1]
  {x * 1.5708};
\addlegendentry{$\theta = (g/g_{\max}) \times \pi/2$ (linear)}

% Região de alta resolução
\fill[red!8, opacity=0.5]
  (axis cs:0,0) rectangle (axis cs:0.3,1.75);

% Anotação θ_Miguel
\addplot[thick, orange, dashed] coordinates
  {(0, 0.1098) (0.1098, 0.1098)};
\addplot[thick, orange, dashed] coordinates
  {(0.1098, 0) (0.1098, 0.1098)};
\fill[orange] (axis cs:0.1098,0.1098) circle (3pt);
\node[font=\scriptsize, orange, anchor=south west]
  at (axis cs:0.13, 0.12)
  {$\thetaM \approx 6{,}29^\circ$};
\node[font=\scriptsize, orange, anchor=north]
  at (axis cs:0.1098, -0.05) {$\sqrt{\alphaii}$};

% Anotação resolução
\node[draw, rounded corners, fill=red!8,
  font=\scriptsize, text width=3cm, align=center]
  at (axis cs:0.15, 1.3)
  {ALTA\\RESOLUÇÃO\\(valores pequenos\\de $g$)};

\node[draw, rounded corners, fill=blue!8,
  font=\scriptsize, text width=3cm, align=center]
  at (axis cs:0.8, 0.6)
  {Menor resolução\\(valores grandes\\de $g$)};

\end{axis}
\end{tikzpicture}
\caption{Função de codificação angular
$\theta = \arcsin(g/g_{\max})$ (vermelho) comparada com
codificação linear (azul tracejado). A região sombreada indica
onde a codificação angular fornece resolução significativamente
superior. O ponto $\thetaM \approx 6{,}29^\circ$ marca a
transição ruído/informação (Constante de Miguel).}
\label{fig:angular}
\end{figure}

% ======================= FIGURA 6 =======================
\begin{figure}[H]
\centering
\begin{tikzpicture}
\begin{axis}[
  width=13cm, height=8cm,
  ybar,
  bar width=12pt,
  ylabel={Taxa de compressão ($\times$)},
  xlabel={Tipo de dado},
  ymin=0, ymax=8.5,
  symbolic x coords={Embeddings, Pesos NN, KV Cache,
    Esparso 95\%, Gradientes},
  xtick=data,
  xticklabel style={font=\small, rotate=15, anchor=east},
  legend style={at={(0.02,0.98)}, anchor=north west,
    font=\small},
  grid=major,
  grid style={gray!15},
  nodes near coords,
  every node near coord/.append style={font=\tiny},
  enlarge x limits=0.15,
]

% v16.1 (linear)
\addplot[fill=blue!40, draw=blue!60]
  coordinates {
    (Embeddings, 3.20)
    (Pesos NN, 2.90)
    (KV Cache, 4.10)
    (Esparso 95\%, 5.80)
    (Gradientes, 3.50)
  };
\addlegendentry{ACOM Linear v16.1}

% v17.0 (Mirror)
\addplot[fill=red!40, draw=red!60]
  coordinates {
    (Embeddings, 3.85)
    (Pesos NN, 3.42)
    (KV Cache, 4.92)
    (Esparso 95\%, 7.14)
    (Gradientes, 4.38)
  };
\addlegendentry{ACOM Mirror v17.0}

\end{axis}

% Anotações de melhoria
\node[font=\scriptsize\bfseries, green!50!black]
  at (2.7, 6.2) {+20,3\%};
\node[font=\scriptsize\bfseries, green!50!black]
  at (5.0, 5.7) {+17,9\%};
\node[font=\scriptsize\bfseries, green!50!black]
  at (7.3, 6.7) {+20,0\%};
\node[font=\scriptsize\bfseries, green!50!black]
  at (9.6, 7.8) {+23,1\%};
\node[font=\scriptsize\bfseries, green!50!black]
  at (11.9, 6.4) {+25,1\%};

% Nota de correlação
\node[draw, rounded corners, fill=yellow!10,
  font=\scriptsize, text width=5cm, align=center]
  at (6.5, 1.5)
  {Correlação $> 0{,}999$ em todos os cenários\\
  Melhoria média: $+21{,}3\%$};

\end{tikzpicture}
\caption{Comparativo de taxa de compressão: ACOM Linear~v16.1
(azul) versus ACOM Mirror~v17.0 (vermelho) para cinco tipos de
dados. A codificação angular supera consistentemente a
quantização linear, com melhoria média de $+21{,}3\%$ e
correlação acima de $0{,}999$.}
\label{fig:comparativo}
\end{figure}

% ======================= FIGURA 7 =======================
\begin{figure}[H]
\centering
\begin{tikzpicture}[
  arrow/.style={-{Stealth[length=3mm]}, very thick},
]

% --- BULK (3D) esquerda ---
\fill[blue!8] (-6.5,-2.5) rectangle (-2,2.5);
\draw[very thick, blue!50] (-6.5,-2.5) rectangle (-2,2.5);
\node[font=\large\bfseries, blue!60] at (-4.25,3)
  {BULK (3D)};
\node[font=\normalsize\bfseries, blue!60] at (-4.25,2)
  {Dado $L$};

% Ondas representando dados 3D
\draw[thick, blue!40, smooth]
  plot[domain=-6:-2.5, samples=50]
  (\x, {0.8*sin(3*\x r) + 0.4*cos(5*\x r)});
\draw[thick, blue!30, smooth]
  plot[domain=-6:-2.5, samples=50]
  (\x, {0.5*sin(4*\x r) - 0.3*cos(2*\x r) - 1});

% Rótulos
\node[font=\scriptsize, blue!50] at (-4.25,-2)
  {Tensor multidimensional};

% --- ESPELHO (centro) ---
\fill[gray!20] (-0.3,-3.5) rectangle (0.3,3.5);
\draw[very thick, gray!60] (-0.3,-3.5) rectangle (0.3,3.5);

% Reflexo do espelho (hachuras)
\foreach \y in {-3,-2.5,...,3}
  \draw[gray!30, thin] (-0.3,\y) -- (0.3,\y+0.3);

\node[font=\normalsize\bfseries, gray!60, rotate=90]
  at (0,0) {ESPELHO};

% --- BOUNDARY (2D) direita ---
\fill[red!8] (2,-2.5) rectangle (6.5,2.5);
\draw[very thick, red!50] (2,-2.5) rectangle (6.5,2.5);
\node[font=\large\bfseries, red!60] at (4.25,3)
  {BOUNDARY (2D)};
\node[font=\normalsize\bfseries, red!60] at (4.25,2)
  {Reflexo $(\psi, \theta)$};

% Estados e ângulos representados
\node[draw, rounded corners, fill=purple!10,
  font=\scriptsize, minimum width=2cm] (psi_box)
  at (4.25,0.5) {$\psi \in \{0,1,2,3\}$};
\node[draw, rounded corners, fill=orange!10,
  font=\scriptsize, minimum width=2cm] (theta_box)
  at (4.25,-0.5) {$\theta \in [0, \pi/2]$};
\node[font=\scriptsize, red!50] at (4.25,-1.5)
  {2\,bits + $b$\,bits/elemento};
\node[font=\scriptsize, red!50] at (4.25,-2)
  {(comprimido)};

% --- SETA REFLECT (Bulk → Boundary) ---
\draw[arrow, blue!60] (-2,1.5) -- (-0.3,1.5)
  node[midway, above, font=\small\bfseries] {REFLECT};
\draw[arrow, blue!60] (0.3,1.5) -- (2,1.5);
\node[font=\scriptsize] at (0,2.3)
  {$g = \sqrt{|L|}$, $\theta = \arcsin(g/g_{\max})$};

% --- SETA MANIFEST (Boundary → Bulk) ---
\draw[arrow, red!60] (2,-1.5) -- (0.3,-1.5)
  node[midway, below, font=\small\bfseries] {MANIFEST};
\draw[arrow, red!60] (-0.3,-1.5) -- (-2,-1.5);
\node[font=\scriptsize] at (0,-2.8)
  {Dobra $\times 2$: $L' = s \times [g_{\max}\sin\theta]^2$};

% --- BULK reconstruído (esquerda, abaixo) ---
\node[draw, dashed, rounded corners, fill=green!8,
  font=\scriptsize, text width=3.5cm, align=center]
  at (-4.25,-3.5)
  {$L' \approx L$\\correlação $\eta > 0{,}999$\\
  (dado re-emerge)};

% Nota conceitual
\node[draw, thick, rounded corners, fill=yellow!10,
  font=\small, text width=6cm, align=center]
  at (0,-4.8)
  {O dado \textbf{não viaja} --- o \textbf{reflexo} viaja.\\
  O dado \textbf{re-emerge} através da\\
  \textbf{Dobra Dimensional} ($\times 2$).};

\end{tikzpicture}
\caption{Diagrama da Dobra Dimensional ($\times 2$): reflexão
especular Bulk$\to$Boundary. O dado $L$ (3D) é projetado no
espelho como reflexo $(\psi, \theta)$ (2D) pela operação
REFLECT, e re-emerge como $L'$ pela operação MANIFEST (Dobra
$\times 2$). Analogia com correspondência holográfica AdS/CFT.}
\label{fig:dobra}
\end{figure}

\newpage
% ============================================================================
\section{Reivindicações}
% ============================================================================

\begin{enumerate}

\item \textbf{MÉTODO DE REFLEXÃO DIMENSIONAL DE DADOS}
compreendendo as etapas de: receber dados de entrada
representando informação a ser processada em formato tensorial;
computar componente gravitacional $g = \sqrt{|L|}$ para cada
elemento $L$ dos dados de entrada; computar derivada discreta
$\partial L$ entre elementos adjacentes; nominar cada elemento
em um estado psiônico~$\psi$ no Espaço de Hilbert
tetradimensional~$\Hfour$ a partir do par
$(\text{sign}(L), \text{sign}(\partial L))$; computar ponto
angular $\theta = \arcsin(g/g_{\max})$, onde $g_{\max}$ é o
máximo global do componente gravitacional; quantizar o ponto
angular~$\theta$ utilizando número predeterminado de bits
angulares; empacotar estados psiônicos e ângulos quantizados em
formato binário; e gerar objeto de reflexão contendo os
componentes empacotados e metadados de reconstrução; o método
\textbf{caracterizado por} a codificação angular utilizar a
função arco-seno $\theta = \arcsin(g/g_{\max})$, onde $\theta$
pertence ao intervalo $[0, \pi/2]$, produzindo resolução de
quantização não-linear com maior precisão para valores pequenos
de~$g$, e pela nominação de estados~$\psi$ no Espaço de
Hilbert~$\Hfour$ compreendendo quatro estados ortonormais com
paridades definidas, e pela reconstrução dos dados através de
desdobramento dimensional
$L' = \text{sign}(L) \times [g_{\max} \times \sin(\theta)]^2$.

\item Método, de acordo com a reivindicação~1,
\textbf{caracterizado pelo fato de} o Espaço de
Hilbert~$\Hfour$ compreender quatro estados psiônicos
ortonormais: estado COLAPSO$^+$~($\psi_0$), correspondente a
$\text{sign}(L) = +1$ e $\text{sign}(\partial L) = -1$, com
vetor base $|0\rangle = (1,0,0,0)$ e paridade inversa~($-1$);
estado ASCENSÃO$^+$~($\psi_1$), correspondente a
$\text{sign}(L) = +1$ e $\text{sign}(\partial L) = +1$, com
vetor base $|1\rangle = (0,1,0,0)$ e paridade normal~($+1$);
estado EMERGÊNCIA$^-$~($\psi_2$), correspondente a
$\text{sign}(L) = -1$ e $\text{sign}(\partial L) = +1$, com
vetor base $|2\rangle = (0,0,1,0)$ e paridade normal~($+1$); e
estado QUEDA$^-$~($\psi_3$), correspondente a
$\text{sign}(L) = -1$ e $\text{sign}(\partial L) = -1$, com
vetor base $|3\rangle = (0,0,0,1)$ e paridade inversa~($-1$).

\item Método, de acordo com a reivindicação~1,
\textbf{caracterizado pelo fato de} a nominação de estados
utilizar construção de índice binário onde
$\text{bit\_high} = 1$ se $L < 0$ e $0$ caso contrário,
$\text{bit\_low} = 1$ se $\partial L > 0$ e $0$ caso contrário,
e o estado $\psi = \text{bit\_high} \times 2 + \text{bit\_low}$,
requerendo exatamente 2~bits por elemento e permitindo
empacotamento de 4~estados por byte.

\item Método, de acordo com a reivindicação~1,
\textbf{caracterizado pelo fato de} ainda compreender o cálculo
da Força de Expulsão $\Fexp$ como média aritmética das
paridades de todos os estados nominados, onde
$\Fexp = (1/N) \times \sum_i \text{paridade}(\psi_i)$, sendo
$\text{paridade}(\psi_0) = \text{paridade}(\psi_3) = -1$ e
$\text{paridade}(\psi_1) = \text{paridade}(\psi_2) = +1$,
fornecendo métrica intrínseca da dinâmica dos dados sem custo
computacional adicional.

\item Método, de acordo com a reivindicação~1,
\textbf{caracterizado pelo fato de} ainda compreender o cálculo
da entropia de estados
$S_\psi = -\sum_k p_k \times \log_2(p_k)$, onde $p_k$ é a
frequência relativa de cada estado $\psi_k$ para
$k \in \{0, 1, 2, 3\}$, fornecendo medida de complexidade da
forma refletida com máximo teórico de 2,0~bits.

\item Método, de acordo com a reivindicação~1,
\textbf{caracterizado pelo fato de} o número predeterminado de
bits angulares ser configurável entre 4 e 16~bits, com valor
preferencial de 8~bits correspondendo a 256~níveis de
quantização angular que oferecem balanço ótimo entre taxa de
compressão e qualidade de reconstrução.

\item Método, de acordo com a reivindicação~1,
\textbf{caracterizado pelo fato de} a compressão entrópica dos
estados psiônicos e dos ângulos quantizados ser realizada
separadamente utilizando compressor sem perda,
preferencialmente zstd em nível de compressão~22 ou zlib em
nível~9.

\item Método, de acordo com a reivindicação~1,
\textbf{caracterizado pelo fato de} utilizar a constante
$\alphaii = 0{,}012$, denominada Constante de Miguel, como
limiar de estabilidade informacional, e o ponto angular
fundamental $\thetaM = \arcsin(\sqrt{\alphaii}) \approx
6{,}29^\circ$ como referência para transição entre ruído e
informação estruturada.

\item Método, de acordo com a reivindicação~1,
\textbf{caracterizado pelo fato de} ser aplicável a tensores
multidimensionais de qualquer formato, incluindo vetores de
embedding de modelos de linguagem, tensores de cache de atenção
(KV cache), matrizes de pesos de redes neurais, dados esparsos
e vetores de gradientes, com preservação de correlação superior
a $0{,}999$ entre dados originais e reconstruídos.

\item \textbf{SISTEMA DE REFLEXÃO DIMENSIONAL DE DADOS}
compreendendo: um módulo de entrada configurado para receber
dados em formato tensorial; um computador gravitacional acoplado
ao módulo de entrada para executar extração $g = \sqrt{|L|}$ e
derivada temporal $\partial L$; um nominador de estados acoplado
ao computador gravitacional para mapear cada par
$(\text{sign}(L), \text{sign}(\partial L))$ em estado
psiônico~$\psi$ no Espaço de Hilbert
tetradimensional~$\Hfour$; um codificador angular acoplado ao
computador gravitacional para computar
$\theta = \arcsin(g/g_{\max})$ e quantizar o ângulo; um
analisador de paridade acoplado ao nominador de estados para
calcular Força de Expulsão $\Fexp$ e Entropia de
Estados~$S_\psi$; um empacotador para serializar estados e
ângulos em formato binário otimizado; um compressor entrópico
para aplicar compressão sem perda; e um módulo de saída para
gerar objeto de reflexão com metadados; o sistema
\textbf{caracterizado por} o codificador angular implementar a
operação $\theta = \arcsin(g/g_{\max})$ que produz codificação
não-linear com maior resolução para valores pequenos de~$g$, e
o nominador de estados implementar mapeamento em Espaço de
Hilbert~$\Hfour$ com quatro estados ortonormais dotados de
paridades definidas.

\item Sistema, de acordo com a reivindicação~10,
\textbf{caracterizado por} ainda compreender um módulo de
manifestação configurado para reconstruir dados originais
através de desdobramento dimensional, compreendendo:
dequantização angular
$\theta = \theta_q / (2^b - 1) \times (\pi/2)$; reconstrução
gravitacional $g = g_{\max} \times \sin(\theta)$; recuperação
de sinal $\text{sign}(L)$ a partir do estado~$\psi$; e dobra
dimensional $L' = \text{sign}(L) \times g^2$.

\item Sistema, de acordo com a reivindicação~10,
\textbf{caracterizado pelo fato de} a implementação
computacional suportar execução tanto em unidade central de
processamento (CPU) via operações vetorizadas NumPy quanto em
unidade de processamento gráfico (GPU) via framework PyTorch,
produzindo resultados numericamente equivalentes.

\item Sistema, de acordo com a reivindicação~10,
\textbf{caracterizado pelo fato de} o objeto de reflexão gerado
pelo módulo de saída compreender: bloco de estados psiônicos
comprimidos, bloco de ângulos quantizados comprimidos, e
metadados incluindo versão do protocolo, forma original do
tensor, número de elementos, bits angulares, valores $g_{\max}$
e $g_{\min}$, Força de Expulsão $\Fexp$, contagem de estados
por tipo, entropia de estados $S_\psi$, e constante~$\alphaii$.

\end{enumerate}

\newpage
% ============================================================================
\section{Resumo}
% ============================================================================

\begin{center}
{\large\bfseries MÉTODO E SISTEMA DE REFLEXÃO DIMENSIONAL\\
DE DADOS BASEADO EM CODIFICAÇÃO ANGULAR\\
NO ESPAÇO DE HILBERT}
\end{center}

\vspace{0.5cm}

A presente invenção refere-se a métodos e sistemas para reflexão
dimensional de dados, aplicáveis aos setores de processamento de
dados, inteligência artificial e compressão de sinais. A
invenção resolve o problema técnico de codificação não-linear
auto-adaptativa de dados, diferenciando-se do estado da técnica
que emprega quantização linear uniforme. O método compreende
interpretar dados de entrada como luz~($L$), extrair componente
gravitacional $g = \sqrt{|L|}$ e derivada discreta~$\partial L$,
nominar cada elemento em estado psiônico~$\psi$ no Espaço de
Hilbert tetradimensional~$\Hfour$ compreendendo quatro estados
ortonormais (COLAPSO$^+$, ASCENSÃO$^+$, EMERGÊNCIA$^-$,
QUEDA$^-$) com paridades definidas, codificar a magnitude via
ponto angular $\theta = \arcsin(g/g_{\max})$ que produz
resolução não-linear com maior precisão para valores pequenos, e
permitir reconstrução por desdobramento dimensional
$L' = \text{sign}(L) \times [g_{\max} \times \sin(\theta)]^2$.
O sistema extrai automaticamente Força de Expulsão~($\Fexp$) e
Entropia de Estados~($S_\psi$) como métricas intrínsecas dos
dados. Testes demonstram superioridade de $21{,}3\%$ em taxa de
compressão sobre quantização linear, com correlação acima de
$0{,}999$, sendo particularmente eficaz para embeddings de
modelos de linguagem, cache de atenção (KV cache) e dados
esparsos em aplicações de inteligência artificial.

\vspace{0.5cm}
\textbf{Figura principal: Figura~\ref{fig:dobra}}
(Diagrama da Dobra Dimensional)

\end{document}
